\section*{章末問題}
\begin{enumerate}[label=\textbf{\thechapter.\arabic*},leftmargin=*]
 \item\label{問:極限と漸化式} 次の極限を求めよ。また，$a_1=1$，$a_{n+1}=\sqrt{2+a_n}$ で定まる数列が単調増加かつ上に有界であることを示し，その極限を求めよ。
 \[
  \text{(1)}\quad\lim_{n\to\infty}\frac{2n^2+3}{5n^2-n},
  \qquad
  \text{(2)}\quad\lim_{n\to\infty}\bigl(\sqrt{n^2+2n}-n\bigr).
 \]
 \item\label{問:等比誤差} $\sum_{k=0}^{\infty}\left(-\frac12\right)^k$ の和を求めよ。$k=0$ から $n-1$ までの和で近似するとき，誤差の絶対値を $10^{-3}$ 未満にする最小の自然数 $n$ を求めよ。
 \item\label{問:連続と微分可能} $x\ne0$ で $f(x)=\frac{\sin2x}{x}$，$x=0$ で $f(0)=c$ とする。$f$ が $0$ で連続となる $c$ を求めよ。また，$g(x)=|x-1|$ は $1$ で連続だが微分可能でないことを，差商を使って説明せよ。
 \item\label{問:微分計算} 次の関数を微分せよ。(2) については $x=0$ における接線も求めよ。
 \[
  \text{(1)}\quad\log(1+x^2),\qquad
  \text{(2)}\quad(x^2+1)\e^{-x},\qquad
  \text{(3)}\quad x^{\sin x}\quad(x>0).
 \]
 \item\label{問:平均値評価} 平均値の定理を使い，任意の実数 $a,b$ に対して $|\sin b-\sin a|\le|b-a|$ が成り立つことを示せ。適用する関数と区間を明記せよ。
 \item\label{問:三次関数グラフ} $f(x)=x^3-3x^2$ の増減，極値，凹凸，変曲点を求めよ。さらに，区間 $[-1,3]$ における最大値・最小値と，それらをとる $x$ を求めよ。
 \item\label{問:指数不等式方程式} 微分を用いて，すべての実数 $x$ について $\e^x\ge1+x$ を示し，等号成立条件を求めよ。また，方程式 $\log x+x=2$ は $x>0$ でただ一つの解をもち，その解が $(1,2)$ にあることを示せ。
 \item\label{問:逆媒介陰} 次を求めよ。
 \begin{enumerate}[label=(\arabic*),leftmargin=*]
  \item $f(x)=x^3+x$ の逆関数を $g$ とするとき，$g'(0)$。
  \item $x=t^2$，$y=t^3$ で表される曲線の，$t=2$ における接線。
  \item 曲線 $x^2+xy+y^2=7$ の点 $(1,2)$ における接線。ただし，その点の近くで $y$ は $x$ の微分可能な関数として表されるものとする。
 \end{enumerate}
\end{enumerate}

\smallskip\noindent 解答は\pageref{解答:第2章}ページ。
